\noindent\textbf{Elección del procedimiento.} La energía cinética es $T=mR^2(\dot\theta^2+\Omega^2\sin^2\theta)/2$ y el potencial $U=mgR(1-\cos\theta)$. Euler--Lagrange da $\ddot\theta=\sin\theta(\Omega^2\cos\theta-g/R)$. Un equilibrio exige $\sin\theta=0$ o $\cos\theta=g/(\Omega^2R)$.

\begin{enumerate}

\item Comprobar que el cociente esté entre -1 y 1 antes de tomar arco coseno.

\[\fitmath{\mathrm{cose}=\frac{g}{{W}^{2}\cdot R}}\]

\[\fitmath{\mathrm{cose}=\frac{(9.81)}{{(6)}^{2}\cdot (0.5)}\approx 0.545\,}\]

\item Convertir a grados la raíz fuera del eje.

\[\fitmath{\mathrm{thetae}=\frac{\arccos\left(\mathrm{cose}\right)\cdot 180}{\pi}}\]

\[\fitmath{\mathrm{thetae}=\frac{\arccos\left((0.545)\right)\cdot 180}{\pi}\approx 56.97534\,{}^{\circ}}\]

\item El equilibrio no axial aparece cuando $\Omega^2R\ge g$.

\[\fitmath{\mathrm{Wcrit}=\sqrt{\frac{g}{R}}}\]

\[\fitmath{\mathrm{Wcrit}=\sqrt{\frac{(9.81)}{(0.5)}}\approx 4.429447\,\mathrm{1/s}}\]

\end{enumerate}

\noindent\textbf{Contraste y comparación.} Además de las posiciones inclinadas simétricas, existen los equilibrios $\theta=0$ y $\pi$. Para esta velocidad, el inferior pierde estabilidad y los inclinados son estables en el modelo sin rozamiento.
